\chapter{Background}
\label{sec:background}

\section{Lean 4 and Formal Proof Checking}

Lean 4 is an interactive theorem prover based on dependent type theory \citep{demoura2021lean4}. Under the propositions-as-types interpretation, a theorem is a type and a proof is a term inhabiting that type. A proof assistant verifies a theorem by elaborating user-facing syntax into typed internal terms and checking them with a trusted kernel.

This distinction between surface syntax and kernel terms is central to the design of GammaZero. The LLM
generates surface-level Lean code: tactic scripts, proof blocks, local declarations, and theorem
applications. Lean then elaborates that code into a lower-level expression whose type can be checked by
the kernel. Many generated programs fail before they ever reach kernel checking because parsing,
elaboration, type-class search, implicit argument inference, or tactic execution fails. An ATP (Automated Theorem Proving) system must therefore extract useful signal from partially successful proof attempts. A fragment may be syntactically
valid but ill-typed, locally useful but incomplete, or complete only because it contains a placeholder.
The framework in this thesis keeps those cases separate instead of collapsing them into a single failure
label.

Lean proof scripts commonly use tactic mode:
\begin{leanblock}
theorem and_comm (P Q : Prop) :
  And P Q -> And Q P := by
  intro h
  exact And.intro h.right h.left
\end{leanblock}

Each tactic transforms a proof state containing local hypotheses and a target goal. If all goals are discharged and the resulting proof term type-checks, Lean accepts the theorem. This kernel-checked execution model is crucial for AI-assisted proving: models may hallucinate, but Lean remains the verifier of final correctness.

Lean also provides placeholder mechanisms for incomplete proofs. The command
\texttt{sorry} lets a proof script elaborate by inserting a placeholder proof term
for an unresolved goal. This is useful during development, but it is not an
accepted final proof in this thesis. Lean also has \texttt{admit}, which is treated
here as a temporary placeholder used only inside scaffolds to stand for sibling
subgoals that are not currently being solved. GammaZero may use these placeholders
to inspect partial proof structure, but a theorem is counted as solved only when
the final Lean check contains neither \texttt{sorry} nor \texttt{admit}.

Lean also makes proof search more sensitive to context. A tactic such as \texttt{exact h} is meaningful
only if \texttt{h} is in scope and has a type compatible with the current goal. A rewrite tactic depends on
which simp lemmas and local equalities are available. A proof by cases or induction can change the shape
of the context and create several new branches. These properties make local proof-state serialization
important: the model must see enough context to generate a plausible action, and the search system must
remember enough context to decide whether two states are really interchangeable.

\section{LLM-based Theorem Proving}

Recent theorem-proving systems increasingly use LLMs to generate tactics, proof terms, or whole proofs \citep{yang2023leandojo}. Whole-proof generation is simple but fragile: one local mistake can invalidate the full output. Search-guided proving is more robust because each candidate step can be checked by Lean, and failed branches can be abandoned.

However, tactic-level search remains expensive. The model must select both tactic names and arguments, such as theorem names for \texttt{rw}, \texttt{exact}, or \texttt{apply}. In larger proofs, the key step is often not a local tactic but the discovery of an intermediate step. This motivates methods that combine neural generation with structured search and explicit decomposition.

There is also a mismatch between how mathematicians describe proofs and how a tactic-level search
system explores them. A human proof often says, "It suffices to prove an auxiliary inequality," or "We
first establish a monotonicity step". But tactic search must discover the auxiliary subgoal, prove it, and connect it back to the original theorem.
This can make the branching factor explode. Proof
skeletons are one way to utilize the high-level reasoning: the model proposes the intermediate
step explicitly, and the search system later solves the exposed subgoal.

\section{Search-based Proof Generation}

Search-guided theorem proving treats Lean as an environment. A state is the current proof state, an action is a generated proof step, and the transition is Lean execution. Common search strategies include beam search, best-first search, and Monte Carlo tree search \citep{polu2020gptf,lample2022hypertree,yang2023leandojo}. These methods differ in how they prioritize frontier nodes, but they share the same core difficulty: the branching factor is large and useful reward may appear only after many decisions.

Proof decomposition naturally introduces an AND/OR search structure. OR structure represents alternative proof steps that can be applied to the current proof state. AND structure appears when a decomposition generates multiple subgoals that must all be solved for the proof to succeed. Compared with a purely linear proof trajectory, an AND/OR representation is better suited for reasoning about intermediate steps, subgoal dependencies, and hierarchical proof structure.

\section{Reinforcement Learning and Sparse Feedback}
Reinforcement learning is a promising direction for automated theorem proving because proof generation naturally involves sequential decision making \citep{kaliszyk2018reinforcement,bansal2019holist}. A system repeatedly chooses proof steps, observes the resulting proof states after Lean execution, and eventually succeeds or fails depending on whether the theorem is fully verified.

However, applying reinforcement learning to Lean proof generation remains difficult. Lean provides extremely sparse feedback: a proof is either accepted or rejected, and useful signal may appear only after many proof steps. Long proofs further increase the difficulty because local actions may influence the final outcome only much later in the search process.

Dense reward signals can help guide search and future policy optimization, but they must be designed carefully. A reward based only on syntactic validity may encourage compilable but logically irrelevant proof fragments. A useful reward mechanism should therefore reflect not only local validity, but also whether generated proof fragments contribute to the final proof structure.
